\subsection{A uniform mixed-product seed of type $(6,3)$}
\label{sec:root-mixed-six-three}
\status{Proved. Complete independent obstruction/crystal and
reduction/descent proof and exact-program audits passed.}

\begin{lemma}[A universal mixed $(6,3)$ Cartier witness]
\label{lem:root-mixed-six-three-cartier}
\status{Proved under stated hypotheses}
Let $k$ be algebraically closed of characteristic two, let
$E_{\rm o}/k$ be any ordinary elliptic curve, and let
$E_{\rm s}:w^2+w=z^3$. On $A=E_{\rm o}\times E_{\rm s}$ put
\[
 N=\OO_{E_{\rm o}}(6O)\boxtimes\OO_{E_{\rm s}}(3O).
\]
A geometric generic smooth $C\in|N|$ has genus $19$, $p$-rank
$17$, and relative stable rank $16$. This includes every ordinary
$j$-invariant, not only a transcendental one.
\end{lemma}
\begin{proof}
Over $k$, an ordinary Weierstrass equation can be normalized to
\[
 E_{\rm o}:y^2+xy=x^3+c,\qquad c\in k^*.
\]
Indeed weighted scaling makes $a_1=1$. Translating $x$ by $a_3$
removes the $a_3y$ term; translating $y$ by the resulting
$x$-coefficient removes that coefficient as well. A further
change $y\mapsto y+u x$ with $u^2+u=a_2$ removes the remaining
quadratic coefficient.
Smoothness is equivalent to $c\ne0$ in this model.
Use the invariant differentials $\omega_{\rm o}=dx/x$ and
$\omega_{\rm s}=dz$, and put $t=c^{1/2}$ and $u=y+t$.
The basis of $H^0(E_{\rm o},\OO(6O))$ is
\[
                    f=(1,x,y,x^2,xy,x^3).
\]
For $v=u/x$ the curve equation gives $y=x^2+xv^2$.
Thus $f_i=A_i^2+xB_i^2$, where
\[
 A=(1,0,x,x,u,0),\qquad B=(0,1,v,0,x,x).
\]
The ordinary Cartier multiplication pairing is therefore
\[
 \frac{C(f_if_j\omega_{\rm o})}{\omega_{\rm o}}
                         =A_iA_j+xB_iB_j.
\]
In the indicated basis its entries, written as functions, are
\[
 \beta_{\rm o}=
 \begin{pmatrix}
 1&0&x&x&u&0\\
 0&x&u&0&x^2&x^2\\
 x&u&y&x^2&0&xu\\
 x&0&x^2&x^2&xu&0\\
 u&x^2&0&xu&xy&x^3\\
 0&x^2&xu&0&x^3&x^3
 \end{pmatrix}.
\]
For the supersingular basis $e=(1,z,w)$ the pairing is
\[
 \beta_{\rm s}=
 \begin{pmatrix}0&1&z\\1&0&w\\z&w&0\end{pmatrix}.
\]
Both formulas follow directly by the $p$-basis Cartier rule;
the second uses $w=w^2+z^3$.

Index the tensor basis $f_i e_j$ by $3i+j$, starting at zero.
For a section $s$ let $T_s$ be the matrix of the inverse-Frobenius
semilinear map
\[
 C_s:H^0(A,N)\longrightarrow H^0(A,N),\qquad
 \sigma\longmapsto
 \frac{C(s\sigma\,\omega_{\rm o}\wedge\omega_{\rm s})}
      {\omega_{\rm o}\wedge\omega_{\rm s}}.
\]
The tensor product of the two displayed pairings specifies every
entry of $T_s$; the coefficient of $s$ is taken to its square root.
Its Serre-dual Frobenius-linear map has matrix
$M_s=(T_s^{[2]})^t$.

Choose the same section for every $c\ne0$:
\[
 s=1+z+x(1+z+w)+xyz+x^3(z+w).
\]
Its ordinary invariant derivative is
$\partial_{
\rm o}s=x+xz+xw+x^3w$.
The two sections are independent, and the displayed pairings give
$C_s(s)=C_s(\partial_{\rm o}s)=0$.
Consequently $\rank M_s\le16$.

There is also an exact uniform lower certificate.
In $M_sM_s^{[2]}$, select the ordered rows and columns
\begin{align*}
 I&=(0,7,2,3,8,5,6,9,11,1,10,14,13,15,4,12),\\
 J&=(0,1,2,3,4,5,6,7,9,10,11,12,13,14,15,16).
\end{align*}
Their $16\times16$ determinant, computed over $\mathbf F_2[c]$, is
\[
                  \det (M_sM_s^{[2]})_{I,J}=c^5.
\]
This is a polynomial identity, not a check at finitely many
parameters. The retained search program verifies it by exact binary
polynomial fraction-free elimination. The second verifier derives
all six $p$-basis identities and all $36$ ordinary-pairing identities
from the curve equation, reconstructs the entire tensor matrix,
checks the two kernel vectors, and verifies the same determinant
using independent polynomial-domain matrix arithmetic.
Since $c\ne0$, the first and second Frobenius images both have
dimension $16$. They are therefore stable, so the stable rank is
exactly $16$.

The nonvanishing minor defines a nonempty open in $|N|$.
The line bundle $N$ is very ample, so its smooth-divisor locus is
also nonempty and open. These two opens meet in the irreducible
space $|N|=\mathbf P^{17}$.
Lemma~\ref{lem:DE-A3-mixed-even-bound} gives the matching generic
upper bound $16$; equivalently its two generic Cartier-kernel
sections give this bound on the same coefficient space.
The universal divisor sequence has $H^1(A,N^{-1})=0$, while
$H^2(A,\OO_A)$ has nilpotent Frobenius. Taking semisimple parts
shows that the stable rank of $H^1(C,\OO_C)$ is the sum of the
stable rank $16$ just computed and the ambient rank
$f(A)=1$. Adjunction gives $g(C)=1+N^2/2=19$.
Thus $f(C)=17$ and the relative stable rank is $16$, as asserted.
\end{proof}

\begin{remark}[Exact certificate and a corrected transfer direction]
\label{rem:root-mixed-six-three-certificate}
\status{Proved under stated hypotheses}
The search program \path{code/mixed_6_3_cartier_search.py} retains
all six tested coefficient vectors and their outputs; its sixth
vector gives the universal minor $c^5$ in
$0.0044763490004697815$ seconds.
The independent construction in
\path{code/mixed_6_3_cartier_verify.py} passed in
$0.18401811700459803$ seconds. Its first execution encountered a
polynomial-domain coercion error in the software's ordinary
remainder routine; the complete failed source and output are
retained. Using pseudo-remainder with the monic curve equation
resolved the implementation error without changing any identity.
The two independent complete proof/program audits reran the search
in $0.002776610002911184$ and $0.002867126000637654$ seconds,
and the symbolic verifier in $0.15830573499988532$ and
$0.16291111600003205$ seconds, respectively. Their separate
outputs are retained without changing the author's outputs.

A proposed norm transfer along an ordinary-factor isogeny of
degree $n$ would send type $(6,3)$ to $(6,3n)$, not $(6n,3)$.
We record that proposal as invalid. The latter class is obtained
by an isogeny on the supersingular factor; that isogeny is etale
only when $n$ is odd. This seed by itself therefore addresses
the previously open thin-strip classes $a\equiv2\pmod4$,
subject to the separate nodal-image argument. It makes no claim
for $a$ divisible by four.
\end{remark}
